\documentclass{ctexart}
% ==================== 依赖与宏包 ====================
\usepackage{graphicx} % 图像处理
\usepackage{geometry} % 页面边距
\usepackage{fancyhdr} % 页眉页脚
\usepackage{lastpage} % 总页数计算
\usepackage{xcolor}   % 字体颜色
\usepackage{listings} % 代码块
\usepackage{tikz}     % 绘制流程图
\usepackage{ctex}     % 中文支持
% 引入处理超链接的宏包
\usepackage[colorlinks=true, linkcolor=black, urlcolor=red]{hyperref}
% 引入首行缩进宏包
\usepackage{indentfirst}
\usetikzlibrary{shapes.geometric, arrows.meta, positioning}

% ==================== 格式与排版配置 ====================
\geometry{left=2.54cm, right=2.54cm, top=2.54cm, bottom=2.54cm }
\pagestyle{fancy} 
\fancyhead{} 
\fancyhead[L]{ Python学习笔记 } % 修正了拼写
\fancyfoot[C]{ 第 ~\thepage 页 (共 ~\pageref{LastPage} 页) } 

\title{Python学习笔记}
\author{钱成发}
\date{2026.4.26}

% 中文标题(小标题)格式设置
\ctexset{
  section={format=\Large\bfseries\raggedright},
  section = {
    name = {,，},        
    number = \chinese{section}, 
    aftername = \hspace{0pt}    
  }
}

% ==================== 编译章节 ====================
\includeonly{Chapter/Entry_Level_Content
}


% ==================== 正文 ====================
\begin{document}

%生成标题和目录
\maketitle
\tableofcontents
\newpage

% ==================== 导言内容 ====================
\section*{导言}
这篇文章是关于Python的学习笔记，由本人亲手使用 \LaTeX{} 编写，旨在记录Python学习过程的知识点，方便自己查阅和复习。本文章中的视频教材为B站UP主“Python编程入门-好学星城”的系列课程，链接如下：
\url{https://www.bilibili.com/video/BV1GWofBcEjM?p=25&vd_source=918c8f2c53176f374e8efeae1d09ded8}
本人认为，Python是一个非常优秀的语言，适合于各种用途，如数据处理、机器学习、爬虫等等。因此我认为有必要学好Python，并且记录下学习过程中的知识点，以便于自己查阅和复习。
文章知识点均由自己编写，部分内容或许不太严谨，或者视频中的部分内容我会直接跳过，还望大家能够理解，但希望文章内容对大家有所帮助。
\par
秉承着万物开源的精神，本文的源代码也会开源，所有内容都会放到本人的GitHub仓库中，欢迎大家 fork 和 star，
链接如下：
\url{https://github.com/ethen887/LaTex-Project}。希望读者可以从中获得到一些帮助，如果大家觉得哪里需要改进，欢迎提出issue。仓库将永久
保持开源，并且持续更新本人的一些笔记或者其它文档，欢迎大家关注。
\newpage

% ==================== 正文内容 ====================







% =================== 显示编译章节内容 ====================
\section{Python 入门技能}
\subsection{Python 介绍}
Python 是一种高级编程语言，由 Guido van Rossum 于 1991 年开发。它以其简洁的语法和强大的功能而闻名，广泛应用于数据分析、人工智能、网络开发等领域。Python 的设计哲学强调代码的可读性和简洁性，使得程序员能够更快地编写和维护代码。
Python 支持多种编程范式，包括面向对象编程、函数式编程和过程式编程。它拥有丰富的标准库和第三方库，使得开发者能够轻松地实现各种功能。Python 的跨平台特性使得它可以在不同的操作系统上运行，如 Windows、macOS 和 Linux。
Python 的应用范围非常广泛，从简单的脚本编写到复杂的 web 开velopment 和数据处理。 Python 的社区非常活跃，拥有丰富的资源，如文档、教程、社区论坛和开源项目。如今，Python 已成为全球最受欢迎的编程语言之一，被广泛应用于各个领域的开发和研究中。



\subsection{Python 安装与环境配置}
要安装 Python，您可以访问 Python 的官方网站（https://www.python.org/）并下载适合您系统的版本。安装完成后，请确保 Python 已经成功安装并配置了环境变量。您可以在命令行中输入以下命令来验证 Python 是否安装成功：
\begin{lstlisting}
    python --version    
如果 Python 已经安装成功，则会显示 Python 的版本号。
\end{lstlisting}


\subsection{开发环境的使用}
在 Python 开发中，选择一个合适的开发环境非常重要。常用的 Python 开发环境包括 PyCharm、Visual Studio Code 和 Jupyter Notebook 等。这些开发环境提供了丰富的功能，如代码自动完成、调试工具和版本控制集成等，可以大大提高开发效率。
PyCharm 是 JetBrains 公司开发的一款强大的 Python IDE，提供了丰富的功能和特性。 Visual Studio Code 是微软开发的一款轻量级的代码编辑器，支持多种编程语言，包括 Python。 Jupyter Notebook 是一个基于 web 的交互式开发环境，特别适合数据分析和科学计算。
选择适合自己的开发环境可以根据个人喜好和项目需求来决定。无论选择哪个开发环境，重要的是要熟悉其功能和使用方法，以便能够高效地进行 Python 开发。作者本人使用的是Visual Studio Code，也就是VS Code，个人感觉非常好用，功能强大，界面简洁，支持多种插件，可以满足各种开发需求。对于初学者来说，VS Code 是一个非常不错的选择。



\subsection{Python 基础语法}
Python 的基础语法非常简单，基本由关键字、标识符、表达式、语句和注释组成。以下是一些基本的 Python 语法：
\begin{lstlisting}[language=Python]
# 这是一个注释，不会被执行
print("Hello, World!") # 这是一个打印语句
x = 10 # 这是一个赋值语句
print(x) # 这是一个打印语句
\end{lstlisting}
\par
在 Python 中，缩进非常重要，它用于表示代码块的开始和结束。通常使用四个空格进行缩进。 Python 还支持多行语句，可以使用反斜杠（\\）来连接多行代码。
\begin{lstlisting}[language=Python]
# 这是一个多行语句的例子
print("This is a") \
    "multiline statement."
    print("This is also a") \
    "multiline statement."
    print("This is a") \
    "multiline statement."
\end{lstlisting}
总之，Python 的基础语法非常简单易学，适合初学者入门。通过掌握 Python 的基础语法，您可以开始编写简单的程序，并逐步深入学习更高级的 Python 功能和库。



\subsection{输出函数}
在 Python 中，输出函数用于将信息显示在屏幕上。最常用的输出函数是 \texttt{print()}，它可以接受多个参数，并将它们以空格分隔的方式输出。以下是一些使用 \texttt{print()} 函数的例子：
\begin{lstlisting}[language=Python]
    print("Hello, World!")
    print("Hello, World!", end = "!") # end参数用于指定输出结束后的字符，默认为换行符
    print("Hello", "World!", sep = "-", end = "!") # sep参数用于指定输出多个参数之间的分隔符，默认为空格
\end{lstlisting}
在上述代码中，\texttt{print()} 函数被调用了三次，分别输出了 "Hello, World!"、“Hello，World!" 和 "Hello-World!"。在Python 中，\texttt{print()} 函数的参数可以接受多个参数，并使用逗号进行分隔。



\subsection{输入函数}
在 Python 中，输入函数用于从用户那里获取输入数据。最常用的输入函数是 \texttt{input()}，它会等待用户输入并返回一个字符串。以下是一个使用 \texttt{input()} 函数的例子：
\begin{lstlisting}[language=Python]
    name = input("请输入你的名字: ")
    print("你好，" + name + "!")
\end{lstlisting}
在上述代码中，\texttt{input()} 函数会等待用户输入，并返回一个字符串。然后，\texttt{print()} 函数会打印出 "你好，" 和用户输入的名字，并加上一个感叹号。
\par
需要注意的是，\texttt{input()} 函数返回的值始终是一个字符串。如果您需要将输入转换为其他数据类型（如整数或浮点数），可以使用相应的类型转换函数，如 \texttt{int()} 或 \texttt{float()}。例如：
\begin{lstlisting}[language=Python]    
    age = int(input("请输入你的年龄: "))
    print("你今年 " + str(age) + " 岁了!")
    height = float(input("请输入你的身高: "))
    print("你身高 " + str(height) + " 厘米!")
\end{lstlisting}



\subsection{列表及常见操作}
在 Python 中，列表是一种常用的数据结构，用于存储多个元素。列表的元素可以是任何类型的数据，如整数、字符串、布尔值等。列表的索引从0开始，并且可以通过索引访问列表中的元素。以下一些常见的列表操作：
\begin{lstlisting}[language=Python]
    # 创建一个列表
    my_list = [1, 2, 3, 4, 5]
    
    # 访问列表中的元素
    print(my_list[0]) # 输出: 1
    print(my_list[2]) # 输出: 3
    
    # 修改列表中的元素
    my_list[1] = 10
    print(my_list) # 输出: [1, 10, 3, 4, 5]
    
    # 添加元素到列表末尾
    my_list.append(6)
    print(my_list) # 输出: [1, 10, 3, 4, 5, 6]
    
    # 插入元素到指定位置
    my_list.insert(2, 20)
    print(my_list) # 输出: [1, 10, 20, 3, 4, 5, 6]
    
    # 删除列表中的元素
    del my_list[3]
    print(my_list) # 输出: [1, 10, 20, 4, 5, 6]
    
    # 获取列表的长度
    print(len(my_list)) # 输出: 6
    
    # 列表切片
    print(my_list[1:4]) # 输出: [10, 20, 4]
    print(my_list[:3]) # 输出: [1, 10, 20]
    print(my_list[3:]) # 输出: [4, 5, 6]
    print(my_list[-2:]) # 输出: [5, 6]
    print(my_list[:-2]) # 输出: [1, 10, 20, 4]
    print(my_list[::2]) # 输出: [1, 20, 5]
\end{lstlisting}
总之，列表是 Python 中非常重要的数据结构，掌握列表的常见操作可以提高 Python 的编程效率。


\subsection{元祖及常见操作}
元组是一种数据结构，用于存储多个元素。元组和列表非常相似，但元组中的元素不能被修改。（类似于C语言中const修饰的数组）以下一些常见的元组操作：
\begin{lstlisting}[language=Python]
    # 创建一个元组
    my_tuple = (1, 2, 3, 4, 5)
    
    # 访问元组中的元素
    print(my_tuple[0]) # 输出: 1
    print(my_tuple[2]) # 输出: 3
    
    # 获取元组的长度
    print(len(my_tuple)) # 输出: 5
    
    # 元组切片
    print(my_tuple[1:4]) # 输出: (2, 3, 4) # 包前不包后
    print(my_tuple[:3]) # 输出: (1, 2, 3)
    print(my_tuple[3:]) # 输出: (4, 5)
    print(my_tuple[-2:]) # 输出: (4, 5)
    print(my_tuple[:-2]) # 输出: (1, 2, 3)
    print(my_tuple[::2]) # 输出: (1, 3, 5)
    # 元组不可变，以下代码会引发错误
    # my_tuple[1] = 10 # TypeError: 'tuple' object does not support item assignment
\end{lstlisting}



\subsection{字典及常见操作}
字典是一种数据结构，用于存储键值对。字典的键必须是唯一的，并且不能重复。以下一些常见的字典操作：
\begin{lstlisting}[language=Python]
    # 创建一个字典
    my_dict = {"name": "Alice", "age": 30, "city": "New York"}
    
    # 访问字典中的值
    print(my_dict["name"]) # 输出: Alice
    print(my_dict["age"]) # 输出: 30
    
    # 修改字典中的值
    my_dict["age"] = 31
    print(my_dict) # 输出: {'name': 'Alice', 'age': 31, 'city': 'New York'}
    
    # 添加新的键值对到字典中
    my_dict["country"] = "USA"
    print(my_dict) # 输出: {'name': 'Alice', 'age': 31, 'city': 'New York', 'country': 'USA'}
    
    # 删除字典中的键值对
    del my_dict["city"]
    print(my_dict) # 输出: {'name': 'Alice', 'age': 31, 'country': 'USA'}
    
    # 获取字典的长度
    print(len(my_dict)) # 输出: 3
    
    # 获取字典中的所有键
    print(my_dict.keys()) # 输出: dict_keys(['name', 'age', 'country'])
    
    # 获取字典中的所有值
    print(my_dict.values()) # 输出: dict_values(['Alice', 31, 'USA'])
    
    # 获取字典中的所有键值对
    print(my_dict.items()) # 输出: dict_items([('name', 'Alice'), ('age', 31), ('country', 'USA')])
\end{lstlisting}
总之，字典是 Python 中非常重要的数据结构，掌握字典的常见操作可以提高 Python 的编程效率。



\subsection{集合及常见操作}
集合是一种数据结构，用于存储多个元素，但不允许重复。集合中的元素必须是不可变的，如字符串、整数、元组等。以下一些常见的集合操作：
\begin{lstlisting}[language=Python]
    # 创建一个集合
    my_set = {1, 2, 3, 4, 5}
    
    # 添加元素到集合中
    my_set.add(6)
    print(my_set) # 输出: {1, 2, 3, 4, 5, 6}
    
    # 从集合中删除元素
    my_set.remove(3)
    print(my_set) # 输出: {1, 2, 4, 5, 6}
    
    # 获取集合的长度
    print(len(my_set)) # 输出: 5
\end{lstlisting}
\newpage
注意：集合中的元素不能重复，如果添加重复的元素，则不会报错，但添加的元素会被忽略；此外，集合中的元素是无序的，因此不能通过索引访问元素。
总之，集合区别于上述的数据结构，集合中的元素不能重复，并且是无序的。集合有着很多的应用场景，掌握集合的常见操作可以提高 Python 的编程效率。
\par
以上就是Python中的一些常见数据结构及其操作，读者要注意区分这四种数据结构，对比四者在定义时的不同，防止混淆，同时希望广大读者能够熟练掌握它们的使用方法。不过Python中的数据结构有很多，这里只列举了一些常用的，更多数据结构请自行查阅。



\subsection{Python中if语句的介绍}
if语句是Python中用于条件判断的语句。它允许程序根据不同的条件执行不同的代码块。if语句的基本语法如下：
\begin{lstlisting}[language=Python]
    if condition:
        # 执行的代码块
    elif another_condition:
        # 执行的代码块
    else:
        # 执行的代码块
        # 注意：if语句可以有多个elif条件，但最多只能有一个else条件。
        # 如果没有满足条件的条件，则执行else代码块。
\end{lstlisting}
在上述代码中，condition是一个表达式，如果condition的值为True，则执行第一行代码块，否则执行else代码块。elif语句用于在if条件不满足时，继续判断另一个条件。如果elif条件满足，则执行对应的代码块，否则继续判断下一个elif条件，直到所有条件都不满足时，执行else代码块。
\par
有关if语句的使用，读者需要注意以下几点：
1. if语句中的条件必须是一个表达式，且其值必须为True或False。
2. if语句可以有多个elif条件，但最多只能有一个else条件。
3. 如果没有满足条件的条件，则执行else代码块。
4. if语句的代码块必须缩进，在condition后面必须有冒号，否则Python会报错。
\par



\subsection{Python中循环语句的介绍}
循环语句是Python中用于重复执行代码的语句。Python提供了三种常用的循环语句：for循环、while循环和range()函数。\par
1. for循环：for循环用于遍历一个序列（如列表、元组、字符串等）或一个可迭代对象。它的基本语法如下：
\begin{lstlisting}[language=Python]
    for variable in sequence:
        # 执行的代码块
        # 注意：for循环中的变量会自动迭代，即每次循环都会将序列中的下一个元素赋给变量。
        # 如果序列中没有元素，则循环结束。
        # 注意：for循环中的代码块必须缩进，最后必须有冒号，否则Python会报错。
\end{lstlisting}
2. while循环：while循环用于在满足条件的情况下重复执行代码块。它的基本语法如下：
\begin{lstlisting}[language=Python]
    while condition:
        # 执行的代码块
        # 注意：while循环中的条件必须为True或False。
        # 如果条件为True，则执行代码块，否则结束循环。
        # 注意：while循环中的代码块必须缩进，最后必须有冒号，否则Python会报错。
\end{lstlisting}
3. range()函数：range()函数用于生成一个整数序列，常用于for循环中。它的基本语法如下：
\begin{lstlisting}[language=Python]
    range(start, stop, step)
    # start: 起始值，默认为0
    # stop: 结束值（不包含），依旧遵循包前不包后的原则
    # step: 步长，默认为1
\end{lstlisting}
例如，以下代码使用for循环和range()函数来打印0到9的数字：
\begin{lstlisting}[language=Python]
    for i in range(10):
        print(i)
\end{lstlisting}


\end{document}